\section{Broad phase}
\label{sec:bp}

Of the two phases the broad one asks the cheaper question. It never evaluates
$\Fdiff$ and never sees a tolerance. It turns a mesh into a set of candidate
pairs, meaning pairs of primitives whose swept bounding boxes overlap and which
might therefore touch, and hands them to \cref{sec:np}. It must be
conservative, because a pair it discards is never examined again. On the scenes
in our benchmark it is also the larger share of the host runtime, so it must be
fast.

One algorithm serves it, over one acceleration structure: a uniform grid on two
of the three axes. The two queries it answers are not the same shape, so they
are described separately. Vertex-face asks one list against another and is
answered by binning a box into every cell its extent touches
(\cref{sec:bp:cell}); edge-edge asks one list against itself, which admits a
cheaper binning and a different walk (\cref{sec:bp:self}).

\subsection{Swept boxes, and the shape of the output}
\label{sec:bp:common}

Because vertices move affinely, the axis-aligned box enclosing a primitive's
entire trajectory is the componentwise hull of its vertex positions at $t=0$ and
$t=1$. That hull is exact and needs no sampling. One box is built per vertex,
per edge and per face. Where the box is stored in a narrower type than the
geometry, each bound is moved outward by one unit in the last place, so a
rounded box can never exclude a candidate the exact box would have kept.

Two conventions shape everything downstream. The first concerns the overlap
test, which rejects only on a \emph{strict} separation, $a_{\min} > b_{\max}$ on
some axis. Boxes that merely touch are therefore reported as overlapping. This
inclusivity is deliberate: a candidate window that rejects a pair of boxes
merely touching drops the pairs a moving surface produces most often, which are
the degenerate ones.

The second is that all storage is structure-of-arrays and every output is built
\emph{count-then-fill}. A first pass counts the pairs each primitive will
contribute, a prefix sum turns those counts into offsets, and a second pass
re-runs the identical traversal, writing into its own slice. Nothing grows a
shared buffer, the allocation is exact, and the output is bit-for-bit
reproducible. The same code shape then works on a GPU, where a growable
per-thread container is not available.

The host tests one box against $32$ consecutive boxes at a time, and an AVX-512,
AVX2 or NEON kernel returns a $32$-bit mask of the survivors. Counting is then a
population count, and enumeration is a bit scan. A rejected candidate therefore
never issues the indirect load of its element indices, which is the expensive
part of a test that usually fails.

\subsection{The vertex-face query}
\label{sec:bp:cell}

Vertex-face takes two lists, the faces against the vertices, and
\cref{alg:cell} answers it by binning one of them into a grid. Two of its
choices are deliberate enough to state before the algorithm itself.

The grid is over \emph{two} axes, not three. For a surface of $N$ elements of
size $h$ inside a region of side $L$ we have $N \sim (L/h)^{2}$, because a
surface does not fill a volume. A two-dimensional grid at element resolution
therefore has $O(N)$ cells, and it is mostly occupied. A three-dimensional one
has $(L/h)^{3} = N^{3/2}$ cells. At $N = 1.5\times10^{6}$ that is
$1.8\times10^{9}$ cells at $0.08\%$ occupancy, which buys memory and empty-cell
iteration in exchange for pruning along a direction the surface barely occupies.
The third axis is not abandoned. The ordinary three-dimensional box test inside
the cell culls it, where it costs nothing extra.

The cell size is derived, not chosen. The grid spacing on each of the two axes
is the \emph{mean extent of the swept boxes} on that axis. The resulting cell
counts are then scaled isotropically so that their product does not exceed $4N$:
past that point more resolution does not pay, and a degenerate extent could
otherwise allocate without bound. No voxel size is left for a caller to supply,
and none to get wrong.

Two other spacings suggest themselves and both are worse, for opposite reasons.
Sizing on the \emph{element} size makes the grid far too fine: a swept box is
much larger than the element it came from, so every box would span a block of
cells and the binning pass would cost more than the query saves. Sizing on the
\emph{largest} box is the textbook uniform-grid rule, and it is the more
tempting of the two, because it buys a property the mean does not -- if no box
is wider than a cell, then two overlapping boxes lie within one cell of each
other and a fixed $3\times3$ stencil around a single owning cell finds every
partner, with no duplicates to settle. What that rule costs is that the widest
box in the scene sets the spacing for all of them. The swept boxes are not of
one size: the widest is $4.5\times$ the mean on armadillo-rollers' vertices and
$24\times$ on cloth-funnel's edges, because a few elements move fast while most
barely move. On armadillo-rollers one swept edge spans $11\%$ of the scene, so a
cell at least that wide admits an $8\times9$ grid at any element count, and the
scene's $36{,}363$ edges land in $72$ cells, $505$ to a cell. Measured, that
costs the host broad phase up to $11.9\times$. The mean keeps the cost of a box
proportional to that box's own size, which is what a heavy-tailed size
distribution requires.

\begin{algorithm}[htbp]
\caption{The vertex-face query, one pass (count or fill)}
\label{alg:cell}
\begin{algorithmic}[1]
\Function{CellList}{$A$, $B$}
  \State $(a_{0},a_{1}) \gets$ the two axes of largest centre variance
  \State $h_{k} \gets$ mean extent of the boxes on axis $a_{k}$;\quad
         $n_{k} \gets \lfloor \operatorname{span}_{k}/h_{k} \rfloor$,
         scaled so $n_{0}n_{1} \le 4N$
  \State count the cells of each box of $B$, prefix-sum them, and scatter
         each box index into its cells
  \ForAll{$i$ \textbf{in parallel}}
    \ForAll{cells $c$ touched by $A_{i}$, and $j$ listed in $c$}
      \If{$A_{i}$ and $B_{j}$ are disjoint in 3D} \textbf{continue} \EndIf
      \State $o_{k} \gets \max(A_{i}.\min_{a_{k}},\, B_{j}.\min_{a_{k}})$
        \Comment{corner of the overlap}
      \If{$c \ne \operatorname{cell}(o_{0},o_{1})$} \textbf{continue}
        \Comment{another cell owns this pair}
      \EndIf
      \If{$A_{i}$ and $B_{j}$ share a vertex} \textbf{continue} \EndIf
      \State count, or emit the pair
    \EndFor
  \EndFor
\EndFunction
\end{algorithmic}
\end{algorithm}

\begin{figure}[htbp]
  \centering
  % The cell list in detail: (a) how the structure is built, (b) how a pair met
% in several cells is emitted once.
%
% No prose inside the drawing. What the panels mean belongs in the caption,
% where it can be read at caption size and does not compete with the marks.
%
% Grids are drawn line by line rather than with TikZ's `grid`, which places its
% lines at multiples of the step measured from the origin and not from the
% corner of the rectangle given to it.
\begin{tikzpicture}[
  font=\scriptsize,
  >={Stealth[length=1.6mm]},
  lbl/.style={font=\scriptsize, text=sccdInk, inner sep=1.5pt},
  tiny/.style={font=\tiny, text=sccdInk, inner sep=1pt},
  ttl/.style={font=\scriptsize\bfseries, text=sccdInk, inner sep=0pt},
  slot/.style={draw=sccdInk!45, line width=0.3pt, minimum width=3.6mm,
               minimum height=3.6mm, inner sep=0pt, font=\tiny, text=sccdInk},
  lead/.style={draw=sccdInk!55, line width=0.3pt},
]

% ============================== (a) construction =============================
% Three boxes over a 3x3 grid, and the two arrays they bin into. Cells are
% numbered row-major from the bottom left, which is the order cell_of gives.
\begin{scope}
  \node[ttl, anchor=south west] at (0,3.30) {(a) construction};

  \foreach \x in {0,0.70,1.40,2.10}
    \draw[draw=sccdInk!30, line width=0.3pt] (\x,0.90) -- (\x,3.00);
  \foreach \y in {0.90,1.60,2.30,3.00}
    \draw[draw=sccdInk!30, line width=0.3pt] (0,\y) -- (2.10,\y);
  \draw[draw=sccdInk!45, line width=0.4pt] (0,0.90) rectangle (2.10,3.00);

  \draw[draw=sccdTeal, line width=0.9pt] (0.15,1.05) rectangle (0.95,1.50);
  \draw[draw=sccdPlum, line width=0.9pt] (0.85,1.75) rectangle (1.55,2.20);
  \draw[draw=sccdAmber, line width=0.9pt] (1.50,2.42) rectangle (1.95,2.88);
  \node[lbl, text=sccdTeal, anchor=west] at (0.99,1.28) {$a$};
  \node[lbl, text=sccdPlum, anchor=west] at (1.59,1.98) {$b$};
  \node[lbl, text=sccdAmber, anchor=west] at (1.99,2.65) {$c$};

  % The cell numbers go on last, so a box cannot hide one.
  \foreach \c/\cx/\cy in {0/0.04/0.94, 1/0.74/0.94, 2/1.44/0.94,
                          3/0.04/1.64, 4/0.74/1.64, 5/1.44/1.64,
                          6/0.04/2.34, 7/0.74/2.34, 8/1.44/2.34}
    \node[tiny, text=sccdInk!55, anchor=south west, fill=white,
          inner sep=0.5pt] at (\cx,\cy) {\c};

  % --- the two arrays ---
  \node[lbl, anchor=east] at (-0.06,0.44) {\texttt{cellptr}};
  \foreach \i/\v in {0/0, 1/1, 2/2, 3/2, 4/2, 5/3, 6/4, 7/4, 8/4, 9/5}
    \node[slot, anchor=south west] at ({0.02 + \i*0.38},0.26) {\v};

  \node[lbl, anchor=east] at (-0.06,-0.16) {\texttt{cellidx}};
  \foreach \i/\v/\col in {0/$a$/sccdTeal, 1/$a$/sccdTeal, 2/$b$/sccdPlum,
                          3/$b$/sccdPlum, 4/$c$/sccdAmber}
    \node[slot, text=\col, anchor=south west] at ({0.02 + \i*0.38},-0.34) {\v};

  % Which cell each run of cellidx belongs to.
  \foreach \i/\c in {0/0, 1/1, 2/4, 3/5, 4/8}
    \node[tiny, text=sccdInk!55, anchor=north] at ({0.21 + \i*0.38},-0.40) {\c};
  \node[tiny, text=sccdInk!55, anchor=east] at (-0.06,-0.50) {cell};
\end{scope}

% ============================ (b) one pair, one cell =========================
\begin{scope}[xshift=57mm]
  \node[ttl, anchor=south west] at (0,3.30) {(b) one pair, one cell};

  \fill[sccdTeal!12] (0.60,1.05) rectangle (1.80,2.85);

  \foreach \x in {0,0.60,1.20,1.80,2.40,3.00}
    \draw[draw=sccdInk!30, line width=0.3pt] (\x,0.45) -- (\x,3.15);
  \foreach \y in {0.45,1.05,1.65,2.25,2.85}
    \draw[draw=sccdInk!30, line width=0.3pt] (0,\y) -- (3.00,\y);
  \draw[draw=sccdInk!45, line width=0.4pt] (0,0.45) rectangle (3.00,3.15);

  \draw[draw=sccdPlum!55, line width=0.5pt, densely dashed] (1.20,1.65) rectangle (2.40,2.85);

  \fill[sccdAmber!45] (1.30,1.80) rectangle (1.70,2.30);

  \draw[draw=sccdTeal, line width=0.9pt] (0.70,1.20) rectangle (1.70,2.30);
  \draw[draw=sccdPlum, line width=0.9pt] (1.30,1.80) rectangle (2.30,2.80);

  % The two cells the pair is met in go on top of the boxes, or the box edges
  % running beside them leave nothing of either mark to see.
  \draw[draw=sccdInk, line width=1.1pt] (1.20,1.65) rectangle (1.80,2.25);
  \draw[draw=sccdInk!45, line width=0.8pt] (1.32,2.37) -- (1.68,2.73);
  \draw[draw=sccdInk!45, line width=0.8pt] (1.68,2.37) -- (1.32,2.73);
  \fill[sccdInk] (1.30,1.80) circle (1.6pt);

  \node[lbl, text=sccdTeal, anchor=east] at (0.66,1.32) {first};
  \node[lbl, text=sccdPlum, anchor=west] at (2.34,2.70) {second};

  % --- key in one column, so no label can run into the next ---
  \fill[sccdAmber!45] (0.00,0.14) rectangle ++(0.22,0.16);
  \node[lbl, anchor=west] at (0.26,0.22) {overlap of the two boxes};
  \fill[sccdInk] (0.11,-0.08) circle (1.6pt);
  \node[lbl, anchor=west] at (0.26,-0.08) {its minimum corner};
  \draw[draw=sccdInk, line width=1.0pt] (0.00,-0.38) rectangle ++(0.22,0.16);
  \node[lbl, anchor=west] at (0.26,-0.30) {the cell that emits the pair};
  \draw[draw=sccdInk!40, line width=0.7pt] (0.02,-0.60) -- (0.20,-0.44);
  \draw[draw=sccdInk!40, line width=0.7pt] (0.20,-0.60) -- (0.02,-0.44);
  \node[lbl, anchor=west] at (0.26,-0.52) {met again here, and dropped};
\end{scope}
\end{tikzpicture}

  \caption{The cell list, schematically and not to scale.
    \textbf{(a)} Construction is count-then-fill, the same shape as the pair
    collection that follows it. One pass over the boxes counts the cells each
    one's extent covers, a prefix sum turns those counts into \texttt{cellptr},
    and a second pass scatters each box index into each of its cells. The
    result is a CRS structure: \texttt{cellptr[c]} to \texttt{cellptr[c+1]} is
    the run of \texttt{cellidx} holding cell $c$'s boxes.
    \textbf{(b)} A box is binned into \emph{every} cell its extent touches, and
    those same cells are what it reads. So when two boxes both reach more than
    one cell in common -- two, here -- the second box's index is listed in each
    of them, and the first box's walk meets it once per shared cell. The pair is
    emitted from the one holding the minimum corner of their overlap.}
  \label{fig:cell2d}
\end{figure}

Construction is count-then-fill, drawn in \cref{fig:cell2d}(a) and the same
shape as the pair collection that follows it. One pass over the boxes counts the
cells each one covers, a prefix sum over the cells turns those counts into the
row pointer, and a second pass scatters each box index into each of its cells.
Nothing is appended to a growing container, so the allocation is exact and known
before a single index is written, and no thread needs a private buffer sized for
a worst case it will not reach. The two passes share one function for deciding
which cells a box covers, so they cannot disagree about it -- a disagreement
there would corrupt the cell array silently rather than fail. The pair output is
built the same way, which is why the whole broad phase allocates exactly twice
and both sizes are counted rather than guessed.

Duplicates are then settled without per-pair state. A box spans several cells,
so a pair can be discovered in several of them, as in \cref{fig:cell2d}(b). The
usual remedies, a hash set of emitted pairs or a mark array, cost memory
proportional to the number of pairs and are awkward on a GPU. Instead the pair
is \emph{attributed} to the cell containing the minimum corner of the two boxes'
overlap, on line 8 of \cref{alg:cell}. That corner lies inside both boxes, so
both are binned in that cell, and the cell is unique. The rule costs two clamps
per surviving pair and no state at all, which is why the same code runs
unchanged on the device.

Nothing on this path is sorted and nothing assumes an order. The query tests
every entry of a touched cell and never breaks on a first miss, so the sort is
removed and not merely moved. Binning costs
$O\bigl(\sum_{i}\text{cells}(i)\bigr)$ and the query
$O\bigl(\sum_{i}\sum_{c \ni i}\text{occupancy}(c)\bigr)$, with no logarithmic
factor in either. Here $\text{cells}(i)$ is the footprint of box $i$ at the mean
spacing, so it is $O(1)$ in the mean and larger for the few boxes that are
larger: a box pays for its own size and for no other box's.
\Cref{sec:results:scaling} measures the exponents both achieve in practice.

\subsection{The edge-edge query}
\label{sec:bp:self}

Edge-edge is one list against itself, and that admits a cheaper binning than
\cref{sec:bp:cell} uses. A box enters the \emph{single} cell holding its
minimum corner rather than every cell its extent touches, so the cell array
holds one entry per box, a partner is met at most once, and the duplicate rule
of \cref{alg:cell} is not needed. Reading
only the cells that come after a box's own is how linked-cell neighbour search
avoids handling a pair twice, a cell visiting half its
neighbours~\citep{allen2017liquids}. It is sound there because a particle is a
point binned by its centre in a cell at least the interaction radius wide, and
extended boxes binned at a corner are neither.

The obvious reading of "after" fails. If a box reads the cells of its own
footprint, its columns crossed with its rows, the pair is found only when one
box's cell lies in the other's footprint --- and two boxes can overlap with
neither doing so. Put one a column to the right of the other and a row below it:
each footprint then runs up and to the right, away from where the other is
binned, and neither box ever reads the other (\cref{fig:incomparable}).
Componentwise order on cells is \emph{partial}, and the incomparable pairs are
exactly the ones lost. Held to its own footprint the walk drops $3{,}698$ of
$23{,}756$ pairs on one of the test cases and $882$ of $6{,}971$ on a
heavy-tailed size spread.


\begin{figure}[htbp]
  \centering
  % Why the walk cannot be bounded by the querying box's own footprint: two
% overlapping boxes whose minimum-corner cells are incomparable.
%
% Conventions as in tikz/cellmin.tex -- no prose inside the drawing, grids
% drawn line by line, and the geometry is real: cells are 0.62 wide, five by
% three, numbered row-major from the bottom left, and every cell index below
% follows from the coordinates in \boxA and \boxB.
\begin{tikzpicture}[
  scale=1.65,
  font=\scriptsize,
  lbl/.style={font=\scriptsize, text=sccdInk, inner sep=1.5pt},
  tiny/.style={font=\tiny, text=sccdInk, inner sep=1pt},
  ttl/.style={font=\scriptsize\bfseries, text=sccdInk, inner sep=0pt},
  idle/.style={draw=sccdInk!22, line width=0.5pt},
  qbox/.style={draw=sccdTeal, line width=1.0pt},
  pbox/.style={draw=sccdPlum, line width=1.0pt},
  read/.style={fill=sccdTeal!14},
  unread/.style={fill=sccdAmber!22},
]

% A: minimum corner in column 2, row 0.  B: column 1, row 1.
\def\boxA{1.50/0.10/1.90/0.90}
\def\boxB{0.80/0.70/2.20/1.00}

\newcommand{\cellfill}[3]{%
  \fill[#3] ({#1*0.62},{#2*0.62}) rectangle ({(#1+1)*0.62},{(#2+1)*0.62});}

\newcommand{\scenegrid}{%
  \foreach \x in {0,0.62,1.24,1.86,2.48,3.10}
    \draw[draw=sccdInk!30, line width=0.3pt] (\x,0) -- (\x,1.86);
  \foreach \y in {0,0.62,1.24,1.86}
    \draw[draw=sccdInk!30, line width=0.3pt] (0,\y) -- (3.10,\y);
  \draw[draw=sccdInk!45, line width=0.4pt] (0,0) rectangle (3.10,1.86);}

\newcommand{\cellnumbers}{%
  \foreach \r in {0,1,2}
    \foreach \c in {0,1,2,3,4} {
      \pgfmathtruncatemacro{\id}{\r*5+\c}
      \node[tiny, text=sccdInk!45, anchor=south west, fill=white,
            inner sep=0.5pt] at ({\c*0.62+0.03},{\r*0.62+0.03}) {\id};}}

% The two boxes, with the minimum corner each is binned at.
\newcommand{\thetwoboxes}{%
  \foreach \xa/\ya/\xb/\yb in \boxA {
    \draw[qbox] (\xa,\ya) rectangle (\xb,\yb);
    \fill[sccdTeal] (\xa,\ya) circle (1.4pt);}
  \foreach \xa/\ya/\xb/\yb in \boxB {
    \draw[pbox] (\xa,\ya) rectangle (\xb,\yb);
    \fill[sccdPlum] (\xa,\ya) circle (1.4pt);}}

% ========================= (a) the two boxes =================================
\begin{scope}
  \node[ttl, anchor=south west] at (0,2.02) {(a) $A$ and $B$ overlap};

  \scenegrid
  % The overlap itself, so that the reader can see there is a pair to find.
  \fill[sccdInk!12] (1.50,0.70) rectangle (1.90,0.90);
  \thetwoboxes
  \cellnumbers
  \node[lbl, text=sccdTeal, anchor=north west] at (1.92,0.45) {$A$};
  \node[lbl, text=sccdPlum, anchor=south west] at (2.22,1.00) {$B$};
\end{scope}

% ===================== (b) A reads its own footprint =========================
% A is binned in cell 2 and its maximum corner in cell 8, so its footprint is
% columns 2-3 crossed with rows 0-1: cells 2, 3, 7, 8. B is binned in cell 6,
% one column to the left of that, and is never read.
\begin{scope}[shift={(3.72,0)}]
  \node[ttl, anchor=south west] at (0,2.02) {(b) $A$'s footprint misses $B$};

  \cellfill{2}{0}{read}  \cellfill{3}{0}{read}
  \cellfill{2}{1}{read}  \cellfill{3}{1}{read}
  \cellfill{1}{1}{unread}
  \scenegrid
  \thetwoboxes
  \cellnumbers
  \node[lbl, text=sccdTeal, anchor=north west] at (1.92,0.45) {$A$};
  \node[lbl, text=sccdPlum, anchor=south west] at (2.22,1.00) {$B$};
\end{scope}

% ===================== (c) B reads its own footprint =========================
% B is binned in cell 6 and its maximum corner in cell 8, so its footprint is
% columns 1-3 on row 1: cells 6, 7, 8. A is binned in cell 2, one row below,
% and is never read.
\begin{scope}[shift={(0,-2.60)}]
  \node[ttl, anchor=south west] at (0,2.02) {(c) $B$'s footprint misses $A$};

  \cellfill{1}{1}{read}  \cellfill{2}{1}{read}  \cellfill{3}{1}{read}
  \cellfill{2}{0}{unread}
  \scenegrid
  \thetwoboxes
  \cellnumbers
  \node[lbl, text=sccdTeal, anchor=north west] at (1.92,0.45) {$A$};
  \node[lbl, text=sccdPlum, anchor=south west] at (2.22,1.00) {$B$};
\end{scope}

% ===================== (d) the walk that does find it ========================
% A holds the smaller linear index, so it reports the pair. Its own row starts
% at its own cell; the row above starts where the row's prefix maximum says
% something reaches back, which is column 1 -- left of A's own column.
\begin{scope}[shift={(3.72,-2.60)}]
  \node[ttl, anchor=south west] at (0,2.02) {(d) $A$ reads cells $2,3$ then $6,7,8$};

  \cellfill{2}{0}{read}  \cellfill{3}{0}{read}
  \cellfill{1}{1}{read}  \cellfill{2}{1}{read}  \cellfill{3}{1}{read}
  \scenegrid
  \thetwoboxes
  \cellnumbers
  \node[lbl, text=sccdTeal, anchor=north west] at (1.92,0.45) {$A$};
  \node[lbl, text=sccdPlum, anchor=south west] at (2.22,1.00) {$B$};
\end{scope}

% ================================== key ======================================
\begin{scope}[shift={(0,-3.30)}]
  \fill[sccdInk] (0.11,0.22) circle (1.4pt);
  \node[lbl, anchor=west] at (0.26,0.22) {the minimum corner a box is binned at};
  \fill[read] (0.00,-0.08) rectangle ++(0.22,0.16);
  \draw[draw=sccdInk!30, line width=0.3pt] (0.00,-0.08) rectangle ++(0.22,0.16);
  \node[lbl, anchor=west] at (0.26,0.00) {the cells the query reads};
  \fill[unread] (0.00,-0.38) rectangle ++(0.22,0.16);
  \draw[draw=sccdInk!30, line width=0.3pt] (0.00,-0.38) rectangle ++(0.22,0.16);
  \node[lbl, anchor=west] at (0.26,-0.30) {the partner's cell, left unread};
\end{scope}
\end{tikzpicture}

  \caption{Why a query is not bounded by its own footprint, drawn to scale over
    a five-by-three grid of cells $0.62$ wide. \textbf{(a)} $A =
    [1.5,1.9]\times[0.1,0.9]$ is binned in cell $2$ and $B =
    [0.8,2.2]\times[0.7,1.0]$ in cell $6$; they overlap on both axes, shaded.
    $B$ begins a column to the left of $A$ and a row above it, so neither
    minimum-corner cell precedes the other componentwise.
    \textbf{(b)} $A$'s footprint is columns $2$--$3$ crossed with rows $0$--$1$,
    cells $2$, $3$, $7$ and $8$. It covers $B$'s row and not its column.
    \textbf{(c)} $B$'s footprint is columns $1$--$3$ on row $1$, cells $6$, $7$
    and $8$. It covers $A$'s column and not its row. Each misses the other by
    one coordinate, and no rule for deciding which box reports a pair recovers
    one that neither box reads.
    \textbf{(d)} Row-major linear index puts $A$ first, so $A$ reports. It reads
    its own row from its own cell, and the row above from column $1$, which the
    row's prefix maximum names because $B$ reaches back to $A$ from there ---
    left of $A$'s own column, which is the step (b) will not take.}
  \label{fig:incomparable}
\end{figure}

Row-major linear index $L(c_{0},c_{1}) = c_{1}n_{0} + c_{0}$ is \emph{total}, and
over it the walk is complete. For an overlapping pair $j$'s minimum does not pass
$i$'s maximum on either grid axis, so $L(\operatorname{cell}(j.\min)) \le
L(\operatorname{cell}(i.\max))$: the row term settles it when the rows differ and
the column does when they agree. Whichever box has the smaller $L$ therefore
holds the other in its range, and equal $L$ means one cell, where the index
decides --- the only index test in \cref{alg:cellmin}.

Read literally that range is whole rows, which is what the bounds are for. Each
cell carries the largest upper bound of the boxes in it and each row a prefix
maximum of those, so a binary search gives the first column holding anything that
reaches back to the querying box and the columns before it are skipped at once. That is a prefix maximum applied per row of cells, and it is what lets the
binning stand without a bound on how many cells a box may span; a grid that stores a box only at its lower-left
cell otherwise caps the box at $k$ cells and widens every query window by $k$ to
match. A second bound per cell does the same on the other axis.

\begin{algorithm}[htbp]
\caption{The edge-edge query, one pass (count or fill)}
\label{alg:cellmin}
\begin{algorithmic}[1]
\Function{EdgeToEdgeBroadPhase}{$A$}
  \State grid and cell size as in \cref{alg:cell}; bin each $A_{i}$ into
         $\operatorname{cell}(A_{i}.\min)$ alone
  \State $u_{c} \gets \max\{A_{j}.\max_{a_{1}} : j \in c\}$
    \Comment{$-\infty$ where $c$ is empty}
  \State $\mu_{c} \gets \max\{A_{j}.\max_{a_{0}} : j \in c',\; c' \le c \text{ in } \operatorname{row}(c)\}$
    \Comment{running maximum along each row}
  \ForAll{$i$ \textbf{in parallel}}
    \State $(c_{0}^{b},c_{1}^{b}) \gets \operatorname{cell}(A_{i}.\min)$;\quad
           $(c_{0}^{e},c_{1}^{e}) \gets \operatorname{cell}(A_{i}.\max)$
    \ForAll{$j$ listed in $(c_{0}^{b},c_{1}^{b})$ \textbf{with} $j > i$}
      \State \Call{Test}{$i$, $j$} \Comment{its own cell, where the index decides}
    \EndFor
    \For{$c_{1} \gets c_{1}^{b}$ \textbf{to} $c_{1}^{e}$}
      \If{$c_{1} = c_{1}^{b}$} $c_{0} \gets c_{0}^{b} + 1$
        \Comment{its own cell is read above}
      \Else{} $c_{0} \gets$ \Call{LowerBound}{$\mu$ on row $c_{1}$ up to $c_{0}^{e}$,
             $A_{i}.\min_{a_{0}}$}
      \EndIf
      \For{$c_{0}$ \textbf{to} $c_{0}^{e}$, cell $c \gets (c_{0},c_{1})$}
        \If{$u_{c} < A_{i}.\min_{a_{1}}$} \textbf{continue}
          \Comment{the cell is not read}
        \EndIf
        \ForAll{$j$ listed in $c$} \Call{Test}{$i$, $j$}
          \Comment{$c$ is ahead of $i$'s cell, so no index test}
        \EndFor
      \EndFor
    \EndFor
  \EndFor
\EndFunction
\Function{Test}{$i$, $j$}
  \If{$A_{i}$ and $A_{j}$ are disjoint in 3D, or share a vertex} \Return \EndIf
  \State count, or emit the pair
\EndFunction
\end{algorithmic}
\end{algorithm}

\begin{figure}[htbp]
  \centering
  % The edge-edge binning worked through on one scene: (a) the binning, then the
% walk of each box that emits a pair, in the order the cells put them.
%
% No prose inside the drawing. What the panels mean belongs in the caption,
% where it can be read at caption size and does not compete with the marks.
%
% Grids are drawn line by line rather than with TikZ's `grid`, which places its
% lines at multiples of the step measured from the origin and not from the
% corner of the rectangle given to it.
%
% The seven boxes and the grid are real: cells are 0.62 wide, six by four,
% numbered row-major from the bottom left, and every cell index, walk range and
% skipped column below follows from the coordinates in \scenebox.
% Drawn at 0.62 to the cell and scaled up to the text width. Coordinates scale,
% node text and line widths do not, so the panels grow and the labels keep their
% size -- which is why the panel offsets below are coordinates and not lengths.
\begin{tikzpicture}[
  scale=1.22,
  font=\scriptsize,
  lbl/.style={font=\scriptsize, text=sccdInk, inner sep=1.5pt},
  tiny/.style={font=\tiny, text=sccdInk, inner sep=1pt},
  ttl/.style={font=\scriptsize\bfseries, text=sccdInk, inner sep=0pt},
  idle/.style={draw=sccdInk!22, line width=0.5pt},
  query/.style={draw=sccdTeal, line width=1.0pt},
  found/.style={draw=sccdPlum, line width=1.0pt},
  full/.style={fill=sccdTeal!14},
  empty/.style={fill=sccdMoss!30},
  beyond/.style={fill=sccdInk!6},
  % The span one binary search over the row's prefix maximum steps over. The
  % tint is translucent so the cell fills beneath it still read: a searched cell
  % can be covered, empty, or neither.
  search/.style={draw=sccdAmber, line width=0.8pt,
                 dash pattern=on 2pt off 1.4pt,
                 fill=sccdAmber, fill opacity=0.16},
]

% The columns a row's lower-bound search steps over, as c0..c1 on row r.
\newcommand{\searchspan}[3]{%
  \draw[search] ({#1*0.62},{#3*0.62}) rectangle ({(#2+1)*0.62},{(#3+1)*0.62});}

% One cell, addressed by column and row.
\newcommand{\cellfill}[3]{%
  \fill[#3] ({#1*0.62},{#2*0.62}) rectangle ({(#1+1)*0.62},{(#2+1)*0.62});}

% label / x0 / y0 / x1 / y1 -- the minimum corner is (x0,y0).
\def\scenebox{%
  a/0.35/0.30/1.10/1.00, b/1.05/0.20/1.65/0.75, c/0.80/0.85/1.50/1.50,
  d/1.95/0.45/2.60/1.55, e/2.40/1.10/3.10/1.70, f/0.30/1.75/0.90/2.30,
  g/0.75/1.95/1.45/2.40}

% The grid, drawn line by line. #1 is unused; kept as a command for brevity.
\newcommand{\scenegrid}{%
  \foreach \x in {0,0.62,1.24,1.86,2.48,3.10,3.72}
    \draw[draw=sccdInk!30, line width=0.3pt] (\x,0) -- (\x,2.48);
  \foreach \y in {0,0.62,1.24,1.86,2.48}
    \draw[draw=sccdInk!30, line width=0.3pt] (0,\y) -- (3.72,\y);
  \draw[draw=sccdInk!45, line width=0.4pt] (0,0) rectangle (3.72,2.48);}

% Every box in the faint style, then the minimum corner each is binned at.
\newcommand{\sceneidle}{%
  \foreach \n/\xa/\ya/\xb/\yb in \scenebox {
    \draw[idle] (\xa,\ya) rectangle (\xb,\yb);
    \fill[sccdInk!30] (\xa,\ya) circle (1.0pt);}}

% ============================ (a) the binning ================================
\begin{scope}
  \node[ttl, anchor=south west] at (0,2.64) {(a) one cell per box};

  \scenegrid
  \foreach \n/\xa/\ya/\xb/\yb in \scenebox {
    \draw[draw=sccdInk!55, line width=0.7pt] (\xa,\ya) rectangle (\xb,\yb);
    \fill[sccdInk] (\xa,\ya) circle (1.4pt);
    \node[lbl, anchor=south west, inner sep=1pt] at ({\xb+0.02},\yb) {$\n$};}

  % Cell numbers last, so no box can hide one, and in the corner the box labels
  % do not use.
  \foreach \r in {0,1,2,3}
    \foreach \c in {0,1,2,3,4,5} {
      \pgfmathtruncatemacro{\id}{\r*6+\c}
      \node[tiny, text=sccdInk!45, anchor=south west, fill=white,
            inner sep=0.5pt] at ({\c*0.62+0.03},{\r*0.62+0.03}) {\id};}
\end{scope}

% ============================= (b) box a walks ===============================
% a is binned in cell 0 and its maximum corner in cell 7, so it covers cells 0,
% 1, 6 and 7 -- rows 0-1 crossed with columns 0-1. On row 1 the running maximum
% starts the scan at column 1: cell 6 is empty, so nothing there reaches back.
\begin{scope}[shift={(4.32,0)}]
  \node[ttl, anchor=south west] at (0,2.64) {(b) $a$ covers cells $0,1,6,7$};

  \cellfill{0}{0}{full}  \cellfill{1}{0}{full}  \cellfill{1}{1}{full}
  \cellfill{0}{1}{empty}
  \fill[beyond] (1.24,0) rectangle (3.72,1.24);
  \scenegrid
  \sceneidle

  \draw[found] (1.05,0.20) rectangle (1.65,0.75);
  \draw[found] (0.80,0.85) rectangle (1.50,1.50);
  \draw[query] (0.35,0.30) rectangle (1.10,1.00);
  \fill[sccdTeal] (0.35,0.30) circle (1.4pt);
  % Row 1: the search starts the row at column 1, stepping over cell 6.
  \searchspan{0}{0}{1}

  \node[lbl, text=sccdPlum, anchor=south west] at (1.67,0.75) {$b$};
  \node[lbl, text=sccdPlum, anchor=south west] at (1.52,1.50) {$c$};
  \node[lbl, text=sccdTeal, anchor=north east] at (0.33,0.28) {$a$};
\end{scope}

% ============================= (c) box d walks ===============================
% d is binned in cell 3 and its maximum corner in cell 16, so it covers cells 3,
% 4, 9, 10, 15 and 16 -- rows 0-2 crossed with columns 3-4. Row 2 holds only f,
% whose maximum on the first axis is 0.90 and so never reaches d: the running
% maximum returns a column past 4 and the whole row is skipped without being
% read.
\begin{scope}[shift={(0,-3.40)}]
  \node[ttl, anchor=south west] at (0,2.64) {(c) $d$ covers cells $3,4,9,10,15,16$};

  % Only 3 and 9 hold anything: d itself and e. Cells 4 and 10 are read and
  % yield nothing, and 15 and 16 are stepped over by the running maximum.
  \cellfill{3}{0}{full}  \cellfill{3}{1}{full}
  \cellfill{4}{0}{empty} \cellfill{4}{1}{empty}
  \cellfill{3}{2}{empty} \cellfill{4}{2}{empty}
  \fill[beyond] (3.10,0) rectangle (3.72,1.86);
  \scenegrid
  \sceneidle

  \draw[found] (2.40,1.10) rectangle (3.10,1.70);
  \draw[query] (1.95,0.45) rectangle (2.60,1.55);
  \fill[sccdTeal] (1.95,0.45) circle (1.4pt);
  % Row 1: the search steps over columns 0-2, left of $d$'s own columns, where a
  % partner binned lower could still reach back. Row 2: it returns a column past
  % 4 and steps over the whole row.
  \searchspan{0}{2}{1}
  \searchspan{0}{4}{2}

  \node[lbl, text=sccdPlum, anchor=south west] at (3.12,1.70) {$e$};
  \node[lbl, text=sccdTeal, anchor=north east] at (1.93,0.43) {$d$};
\end{scope}

% ============================= (d) box f walks ===============================
% f is binned in cell 12 and its maximum corner in cell 19, so it covers cells
% 12, 13, 18 and 19 -- rows 2-3 crossed with columns 0-1. On row 3 the running
% maximum starts at column 1, cell 18 being empty.
\begin{scope}[shift={(4.32,-3.40)}]
  \node[ttl, anchor=south west] at (0,2.64) {(d) $f$ covers cells $12,13,18,19$};

  \cellfill{0}{2}{full}  \cellfill{1}{3}{full}
  \cellfill{1}{2}{empty} \cellfill{0}{3}{empty}
  \fill[beyond] (1.24,1.24) rectangle (3.72,2.48);
  \scenegrid
  \sceneidle

  \draw[found] (0.75,1.95) rectangle (1.45,2.40);
  \draw[query] (0.30,1.75) rectangle (0.90,2.30);
  \fill[sccdTeal] (0.30,1.75) circle (1.4pt);
  % Row 3: the search starts the row at column 1, stepping over cell 18.
  \searchspan{0}{0}{3}

  \node[lbl, text=sccdPlum, anchor=south west] at (1.47,2.40) {$g$};
  \node[lbl, text=sccdTeal, anchor=north east] at (0.28,1.73) {$f$};
\end{scope}

% ================================== key ======================================
\begin{scope}[shift={(0,-4.30)}]
  \fill[sccdInk] (0.11,0.22) circle (1.4pt);
  \node[lbl, anchor=west] at (0.26,0.22) {the minimum corner a box is binned at};
  \fill[full] (0.00,-0.08) rectangle ++(0.22,0.16);
  \draw[draw=sccdInk!30, line width=0.3pt] (0.00,-0.08) rectangle ++(0.22,0.16);
  \node[lbl, anchor=west] at (0.26,0.00) {a covered cell holding a box};
  \fill[empty] (0.00,-0.38) rectangle ++(0.22,0.16);
  \draw[draw=sccdInk!30, line width=0.3pt] (0.00,-0.38) rectangle ++(0.22,0.16);
  \node[lbl, anchor=west] at (0.26,-0.30) {a covered cell that is empty};
  \fill[beyond] (0.00,-0.68) rectangle ++(0.22,0.16);
  \draw[draw=sccdInk!30, line width=0.3pt] (0.00,-0.68) rectangle ++(0.22,0.16);
  \node[lbl, anchor=west] at (0.26,-0.60) {past the box's own maximum corner};
  \draw[search] (0.00,-0.98) rectangle ++(0.22,0.16);
  \node[lbl, anchor=west] at (0.26,-0.90)
    {the columns one lower-bound search steps over};
\end{scope}
\end{tikzpicture}

  \caption{The edge-edge binning run through on one scene, drawn to scale: every
    cell index, walk range and skipped span below follows from the coordinates
    shown. \textbf{(a)} Seven swept boxes over a six-by-four grid, each binned
    into the one cell holding its minimum corner --- $a$ into cell $0$, $b$ into
    $1$, $d$ into $3$, $c$ into $7$, $e$ into $9$, $f$ into $12$ and $g$ into
    $19$, the cells numbered row-major from the bottom left. Four pairs of these
    boxes overlap, and each is emitted by whichever of the two sits in the lower
    cell, so only three of the seven boxes emit anything.
    \textbf{(b)} $a$ covers cells $0$, $1$, $6$ and $7$, the last holding its
    maximum corner, and finds $b$ and $c$ in $1$ and $7$. Cell $6$ is empty, so
    the running maximum starts that row at the next column and the cell is never
    read. A run of empty cells costs one binary search and not a test per cell.
    \textbf{(c)} $d$ covers cells $3$, $4$, $9$, $10$, $15$ and $16$, of which
    $3$ holds $d$ itself and $9$ holds $e$, the pair it emits. Cells $4$ and
    $10$ are read and yield nothing. The third row is skipped whole, the empty
    cells $15$ and $16$ included although $d$ covers them: the
    only box binned in that row is $f$, which ends well to the left of where $d$
    begins, so one search puts the row past $d$'s own maximum column and it is
    never entered. Past its own row the walk also reaches left of its own
    columns, because a box binned in a lower column can still be wide enough to
    reach back into $d$ and $d$ holds the lower cell, so that pair would be its
    to report; those cells are skipped by the same search and are left unmarked
    here, being too far from $d$ for the drawing to say anything about them.
    \textbf{(d)} $f$ covers cells $12$, $13$, $18$ and $19$, and finds $g$ in
    $19$; of the two empty cells $13$ is read and $18$ the search steps over.
    $b$, $c$, $e$ and $g$ walk too, and emit
    nothing: each of their partners sits in a lower cell and has already
    reported the pair.}
  \label{fig:cellmin}
\end{figure}

Cost follows the same two terms as \cref{alg:cell}, with the binning reduced to
one incidence per box and the query to whatever the bounds leave. Ordering each
cell's entries on the axis the
grid does not use would make the scan inside a cell monotone as well --- a grid
and sweep-and-prune hybrid of the kind \citet{serpa2020broadmark} evaluates. It
culls where cells are densely occupied and pays an ordering pass everywhere, and
on our scenes the two roughly cancel, so the cells are left unordered.
